%% AMS-LaTeX Created with the Wolfram Language : www.wolfram.com

\documentclass[12pt]{article}

\usepackage[letterpaper,margin=1in]{geometry}
\usepackage{amsmath, amssymb, graphics, setspace}
\usepackage{graphicx} 
\graphicspath{ {./figs/} }
\usepackage{tcolorbox}
\usepackage[framed,numbered,autolinebreaks,useliterate]{mcode}
\newcommand{\mathsym}[1]{{}}
\newcommand{\unicode}[1]{{}}

\usepackage{hyperref}
\hypersetup{
    colorlinks=true,
    linkcolor=blue,
    filecolor=magenta,      
    urlcolor=cyan,
    pdftitle={Overleaf Example},
    pdfpagemode=FullScreen,
}

\def\mathbi#1{\textbf{\em #1}}

\begin{document}
\doublespacing	

\title{CE 401/5501 Homework 09\\
\vspace{2in}
\textbf{Name: \rule{2in}{0.15mm}}}
\author{}
\maketitle

\newcommand*{\I}{\imath} % imaginary unit i

\newpage
\noindent Please review lecture notes \href{https://umsystem.instructure.com/courses/280496/files/33335361?module_item_id=9867784}{Topic 7} and complete the following ML model design problems. 

\section*{Problem 1: Multi-class Classification}

\begin{enumerate}
    \item \textbf{One-Hot Encoding:}\\
    We assign each bridge health condition an exclusive one-hot vector:
    \[
    \text{Excellent Condition: } [1, 0, 0], \quad \text{Safe Condition: } [0, 1, 0], \quad \text{Unsafe Condition: } [0, 0, 1].
    \]
    
    \item \textbf{Softmax Calculation:}\\
    Suppose the model outputs scores:
    \[
    z_1 = 10, \quad z_2 = 5, \quad z_3 = -5.
    \]
    The softmax function is defined as:
    \[
    \text{softmax}(z_i) = \frac{e^{z_i}}{\sum_{j=1}^{3} e^{z_j}}.
    \]
    We compute:
    \[
    e^{10} \approx 22026.47,\quad e^{5} \approx 148.41,\quad e^{-5} \approx 0.0067.
    \]
    Therefore, the denominator is:
    \[
    22026.47 + 148.41 + 0.0067 \approx 22174.89.
    \]
    The softmax outputs are:
    \[
    \text{softmax}(z_1) \approx \frac{22026.47}{22174.89} \approx 0.9933,
    \]
    \[
    \text{softmax}(z_2) \approx \frac{148.41}{22174.89} \approx 0.0067,
    \]
    \[
    \text{softmax}(z_3) \approx \frac{0.0067}{22174.89} \approx 3.0\times10^{-7}.
    \]
    
    \item \textbf{Predicted Class:}\\
    Since the highest softmax value is for \(z_1\), the model predicts the bridge is in ``Excellent Condition.''
\end{enumerate}

\newpage

\section*{Problem 2: Analysis of Model Performance}

We test a model on a dataset of 100 data points with the following true distribution:
\[
\text{Excellent Condition (Level 1): }70,\quad \text{Safe Condition (Level 2): }25,\quad \text{Unsafe Condition (Level 3): }5.
\]
Note that in the provided confusion matrices the \textbf{true labels are arranged in the columns} and the predicted labels in the rows.

\subsection*{(1) Validity Check of the Confusion Matrices}

\textbf{ML Model's Confusion Matrix:}
\[
\begin{bmatrix}
60 & 8 & 2 \\
10 & 14 & 1 \\
0 & 3 & 2 \\
\end{bmatrix}
\]
\textbf{Random Model's Confusion Matrix:}
\[
\begin{bmatrix}
24 & 8 & 0 \\
23 & 9 & 1 \\
23 & 8 & 4 \\
\end{bmatrix}
\]

Since the true labels are placed in the columns, we check the column sums:
\begin{itemize}
    \item For the ML model:
    \begin{itemize}
        \item Column 1 (True Excellent): \(60+10+0 = 70\)
        \item Column 2 (True Safe): \(8+14+3 = 25\)
        \item Column 3 (True Unsafe): \(2+1+2 = 5\)
    \end{itemize}
    \item For the random model:
    \begin{itemize}
        \item Column 1 (True Excellent): \(24+23+23 = 70\)
        \item Column 2 (True Safe): \(8+9+8 = 25\)
        \item Column 3 (True Unsafe): \(0+1+4 = 5\)
    \end{itemize}
\end{itemize}
Both confusion matrices correctly reflect the true label distribution.

\subsection*{(2) Overall Accuracy (OA)}

Overall Accuracy is calculated as the ratio of correct predictions (diagonal elements) to the total number of data points.
\[
\text{OA} = \frac{\text{Diagonal Sum}}{100}.
\]
\begin{itemize}
    \item \textbf{ML Model:} Diagonal sum \(= 60 + 14 + 2 = 76\), so
    \[
    \text{OA}_{\text{ML}} = \frac{76}{100} = 76\%.
    \]
    \item \textbf{Random Model:} Diagonal sum \(= 24 + 9 + 4 = 37\), so
    \[
    \text{OA}_{\text{Random}} = \frac{37}{100} = 37\%.
    \]
\end{itemize}
Thus, based solely on overall accuracy, the ML model outperforms the random model.

\subsection*{(3) Performance for \textit{Excellent Condition} (Level 1)}

For class Excellent (Level 1):
\begin{itemize}
    \item \textbf{True Positives (TP)}: Instances predicted as Excellent that are truly Excellent: \(TP = 60\).
    \item \textbf{False Positives (FP)}: Instances predicted as Excellent but whose true label is not Excellent. These are the other entries in row 1:
    \[
    FP = 8+2 = 10.
    \]
    \item \textbf{False Negatives (FN)}: True Excellent instances that were predicted as another class. These are the other entries in column 1:
    \[
    FN = 10+0 = 10.
    \]
\end{itemize}
Therefore,
\[
\text{Precision}_{\text{Excellent}} = \frac{60}{60+10} \approx 0.857 \quad (85.7\%)
\]
\[
\text{Recall}_{\text{Excellent}} = \frac{60}{60+10} \approx 0.857 \quad (85.7\%)
\]

\subsection*{(4) Performance for \textit{Unsafe Condition} (Level 3)}

For class Unsafe (Level 3):
\begin{itemize}
    \item \textbf{True Positives (TP)}: Instances predicted as Unsafe that are truly Unsafe: \(TP = 2\).
    \item \textbf{False Positives (FP)}: Instances predicted as Unsafe when the true label is not Unsafe. These are the other entries in row 3:
    \[
    FP = 0+3 = 3.
    \]
    \item \textbf{False Negatives (FN)}: True Unsafe instances that were predicted as another class. These are the other entries in column 3:
    \[
    FN = 2+1 = 3.
    \]
\end{itemize}
Thus,
\[
\text{Precision}_{\text{Unsafe}} = \frac{2}{2+3} = 0.4 \quad (40\%)
\]
\[
\text{Recall}_{\text{Unsafe}} = \frac{2}{2+3} = 0.4 \quad (40\%)
\]

\subsection*{(5) Evaluation of the Random-Guess Model}

Even though a random-guess model might occasionally show higher recall for a particular class due to chance, note the following:
\begin{itemize}
    \item The confusion matrix for the random model is correctly constructed (true labels in columns), and its column sums match the true label distribution.
    \item The overall accuracy is very low (37\% compared to 76\% for the ML model).
    \item For a uniform random guess among three classes, one would expect similar performance across classes; any anomalously high recall in a minority class (e.g., Unsafe with only 5 instances) would be due to statistical fluctuations given the small sample size.
\end{itemize}
Thus, the random model is not competitive overall, and any seemingly higher recall for a minority class is likely a result of chance rather than superior performance.

\end{document}
